\documentclass[12pt,a4paper]{article}

\usepackage[margin=1in]{geometry}
\usepackage[T1]{fontenc}
\usepackage[utf8]{inputenc}
\usepackage{lmodern}
\usepackage{enumitem}
\usepackage{hyperref}

\title{Functionality (30\%)\\Jobra Hospital Management System}
\author{}
\date{\today}

\begin{document}
\maketitle

\section{Feature completeness}
\begin{itemize}[leftmargin=*, itemsep=0.4em]
  \item \textbf{Multi-role system}: The application supports \textbf{admin}, \textbf{doctor}, and \textbf{patient} roles with separate dashboards and protected navigation.
  \item \textbf{Authentication flow}: Users can register (patients) and log in, receiving a JWT token used for subsequent requests.
  \item \textbf{Appointments}: Patients can browse doctors and book appointments; doctors can view appointments and update status; admins can view and manage appointments system-wide.
  \item \textbf{Medical reports}: Doctors can create/update reports (diagnosis, prescription, tests, advice) and patients can view their reports.
  \item \textbf{Complaints / reviews}: Patients can submit complaints/feedback and view their own submissions; doctors can view reviews about themselves; admins can view all complaints.
  \item \textbf{Profile / personal info}: The frontend includes ``My Info'' and ``Edit Info'' screens for each role, supporting basic profile management.
\end{itemize}

\section{User experience quality}
\begin{itemize}[leftmargin=*, itemsep=0.4em]
  \item \textbf{Single-page experience}: The frontend is implemented as a React SPA using React Router, providing fast navigation without full page reloads.
  \item \textbf{Role-specific layout}: A shared \texttt{DashboardLayout} component gives consistent navigation and structure across admin/doctor/patient pages.
  \item \textbf{Auto-login routing behavior}: If a valid user session exists in \texttt{localStorage}, the app redirects users directly to their role dashboard instead of showing login/register again.
  \item \textbf{Service abstraction}: API calls are encapsulated in \texttt{services/} modules, reducing duplication and improving UI code readability.
\end{itemize}

\section{Error handling}
\begin{itemize}[leftmargin=*, itemsep=0.4em]
  \item \textbf{Backend status codes}: Controllers return meaningful HTTP codes for common cases (e.g., \texttt{400} for missing fields, \texttt{401} for invalid credentials or missing/invalid token, \texttt{403} for forbidden role access, \texttt{500} for database/server errors).
  \item \textbf{Try/catch in async flows}: Several endpoints use \texttt{try/catch} patterns around database operations to return consistent JSON error responses.
  \item \textbf{Improvement note}: Some endpoints return raw DB errors directly. Standardizing error responses (and adding centralized error middleware) would improve consistency and security (avoid leaking SQL details).
\end{itemize}

\section{Cross-browser compatibility}
\begin{itemize}[leftmargin=*, itemsep=0.4em]
  \item \textbf{CRA browserslist}: The React app uses Create React App tooling and includes a \texttt{browserslist} configuration to target modern production browsers.
  \item \textbf{Standards-based stack}: The UI is built with React + CSS and uses widely supported web APIs; Axios handles HTTP requests in a cross-browser way.
  \item \textbf{Practical expectation}: With the current setup, the app should run reliably on recent Chrome/Firefox/Safari/Edge versions; any legacy-browser support would require additional testing and polyfills.
\end{itemize}

\end{document}

